
\documentclass{article}
\usepackage{colortbl}
\usepackage{makecell}
\usepackage{multirow}
\usepackage{supertabular}

\begin{document}

\newcounter{utterance}

\centering \large Interaction Transcript for game `codenames', experiment `risk\_low', episode 1 with gpt{-}5.4{-}azure.
\vspace{24pt}

{ \footnotesize  \setcounter{utterance}{1}
\setlength{\tabcolsep}{0pt}
\begin{supertabular}{c@{$\;$}|p{.15\linewidth}@{}p{.15\linewidth}p{.15\linewidth}p{.15\linewidth}p{.15\linewidth}p{.15\linewidth}}
    \# & $\;$A & \multicolumn{4}{c}{Game Master} & $\;\:$B\\
    \hline

    \theutterance \stepcounter{utterance}  
    & & & \multicolumn{2}{p{0.3\linewidth}}{
        \cellcolor[rgb]{0.9,0.9,0.9}{
            \makecell[{{p{\linewidth}}}]{
                \texttt{\tiny{[GM$|$GM]}}
                {'hidden': {'team': ['брюки', 'суд', 'смотреть', 'солома', 'шип', 'адвокат', 'бумага', 'флейта', 'замок'], 'innocent': ['лук', 'крест', 'стадион', 'крот', 'экран', 'бассейн', 'треугольник', 'пластик'], 'opponent': ['акции', 'ртуть', 'пират', 'труба', 'звук', 'ветер', 'мяч', 'переключатель'], 'assassin': []}, 'revealed': {'team': {'team': [], 'innocent': [], 'opponent': [], 'assassin': []}, 'opponent': {'team': [], 'innocent': [], 'opponent': [], 'assassin': []}}}
            }
        }
    }
    & & \\ \\

    \theutterance \stepcounter{utterance}  
    & & \multicolumn{4}{p{0.6\linewidth}}{
        \cellcolor[rgb]{0.9,0.9,0.9}{
            \makecell[{{p{\linewidth}}}]{
                \texttt{\tiny{[P1$\langle$GM]}}
                \texttt{Дайте одну подсказку, состоящую из одного слова, которая относится к некоторым словам вашей команды из списка ниже. Вы можете выбрать, к скольким и к каким именно словам относится ваша подсказка.} \\
\texttt{Цель — чтобы ваш партнёр угадал все слова вашей команды на основе подсказок, которые вы ему даёте, прежде чем другая команда раскроет все свои. Если ваш партнёр угадает слово 'соперника' или слово{-}'отвлекатор', его ход угадывания заканчивается, и команда соперников получает право сделать ход. Если ваш партнёр угадает слово{-}'убийца', игра заканчивается, и ваша команда немедленно проигрывает.} \\
\texttt{Чем к большему числу слов относится ваша подсказка, тем быстрее вы можете выиграть игру, но тем сложнее вашему партнёру может быть угадать ассоциации, поэтому выбирайте слово подсказки с умом. Слово подсказки должно быть семантически связано с целевыми словами; оно не может быть одним из слов в списках или содержать их части.} \\
\\ 
\texttt{Всегда давайте вашу однословную подсказку и ваш список связанных целевых слов, разделённых запятыми, в следующем формате, и делайте ответы как можно короче; никогда не включайте никакого другого текста, кроме того, что требуется в этой форме:} \\
\\ 
\texttt{ПОДСКАЗКА: <СЛОВО>} \\
\texttt{ЦЕЛИ: <ЦЕЛИ>} \\
\\ 
\texttt{Слова вашей команды: брюки, суд, смотреть, солома, шип, адвокат, бумага, флейта, замок.} \\
\texttt{Слова команды соперника: акции, ртуть, пират, труба, звук, ветер, мяч, переключатель.} \\
\texttt{Слова{-}отвлекаторы: лук, крест, стадион, крот, экран, бассейн, треугольник, пластик.} \\
\texttt{Слова{-}убийцы: .} \\
            }
        }
    }
    & \\ \\

    \theutterance \stepcounter{utterance}  
    & \multicolumn{4}{p{0.6\linewidth}}{
        \cellcolor[rgb]{0.9,0.9,0.9}{
            \makecell[{{p{\linewidth}}}]{
                \texttt{\tiny{[P1$\rangle$GM]}}
                \texttt{ПОДСКАЗКА: юрист} \\
\texttt{ЦЕЛИ: суд, адвокат} \\
            }
        }
    }
    & & \\ \\

    \theutterance \stepcounter{utterance}  
    & & & \multicolumn{2}{p{0.3\linewidth}}{
        \cellcolor[rgb]{0.9,0.9,0.9}{
            \makecell[{{p{\linewidth}}}]{
                \texttt{\tiny{[GM$|$GM]}}
                \texttt{юрист} \\
            }
        }
    }
    & & \\ \\

    \theutterance \stepcounter{utterance}  
    & & & \multicolumn{2}{p{0.3\linewidth}}{
        \cellcolor[rgb]{0.9,0.9,0.9}{
            \makecell[{{p{\linewidth}}}]{
                \texttt{\tiny{[GM$|$GM]}}
                ['суд', 'адвокат']
            }
        }
    }
    & & \\ \\

    \theutterance \stepcounter{utterance}  
    & & & \multicolumn{2}{p{0.3\linewidth}}{
        \cellcolor[rgb]{0.9,0.9,0.9}{
            \makecell[{{p{\linewidth}}}]{
                \texttt{\tiny{[GM$|$GM]}}
                {'word': 'суд', 'assignment': 'team'}
            }
        }
    }
    & & \\ \\

    \theutterance \stepcounter{utterance}  
    & & & \multicolumn{2}{p{0.3\linewidth}}{
        \cellcolor[rgb]{0.9,0.9,0.9}{
            \makecell[{{p{\linewidth}}}]{
                \texttt{\tiny{[GM$|$GM]}}
                {'word': 'адвокат', 'assignment': 'team'}
            }
        }
    }
    & & \\ \\

    \theutterance \stepcounter{utterance}  
    & & \multicolumn{4}{p{0.6\linewidth}}{
        \cellcolor[rgb]{0.9,0.9,0.9}{
            \makecell[{{p{\linewidth}}}]{
                \texttt{\tiny{[GM$\rangle$P2]}}
                \texttt{Предоставьте список из не более чем 2 слов, разделённых запятыми, из следующего списка, которые лучше всего относятся или наиболее тесно связаны со словом «юрист». Всегда начинайте список вашей(их) догадки(ок) с ОТВЕТ:  и не включайте в свой ответ никакого другого текста.} \\
\\ 
\texttt{замок, ветер, шип, пират, крот, переключатель, бассейн, звук, акции, экран, брюки, пластик, солома, ртуть, адвокат, труба, мяч, смотреть, крест, треугольник, лук, бумага, стадион, флейта, суд} \\
            }
        }
    }
    & \\ \\

    \theutterance \stepcounter{utterance}  
    & & & \multicolumn{4}{p{0.6\linewidth}}{
        \cellcolor[rgb]{0.9,0.9,0.9}{
            \makecell[{{p{\linewidth}}}]{
                \texttt{\tiny{[GM$\langle$P2]}}
                \texttt{ОТВЕТ: адвокат, суд} \\
            }
        }
    }
     \\ \\

    \theutterance \stepcounter{utterance}  
    & & & \multicolumn{2}{p{0.3\linewidth}}{
        \cellcolor[rgb]{0.9,0.9,0.9}{
            \makecell[{{p{\linewidth}}}]{
                \texttt{\tiny{[GM$|$GM]}}
                ['адвокат', 'суд']
            }
        }
    }
    & & \\ \\

    \theutterance \stepcounter{utterance}  
    & & & \multicolumn{2}{p{0.3\linewidth}}{
        \cellcolor[rgb]{0.9,0.9,0.9}{
            \makecell[{{p{\linewidth}}}]{
                \texttt{\tiny{[GM$|$GM]}}
                {'word': 'адвокат', 'assignment': 'team'}
            }
        }
    }
    & & \\ \\

    \theutterance \stepcounter{utterance}  
    & & & \multicolumn{2}{p{0.3\linewidth}}{
        \cellcolor[rgb]{0.9,0.9,0.9}{
            \makecell[{{p{\linewidth}}}]{
                \texttt{\tiny{[GM$|$GM]}}
                {'word': 'адвокат', 'assignment': 'team'}
            }
        }
    }
    & & \\ \\

    \theutterance \stepcounter{utterance}  
    & & & \multicolumn{2}{p{0.3\linewidth}}{
        \cellcolor[rgb]{0.9,0.9,0.9}{
            \makecell[{{p{\linewidth}}}]{
                \texttt{\tiny{[GM$|$GM]}}
                {'word': 'суд', 'assignment': 'team'}
            }
        }
    }
    & & \\ \\

    \theutterance \stepcounter{utterance}  
    & & & \multicolumn{2}{p{0.3\linewidth}}{
        \cellcolor[rgb]{0.9,0.9,0.9}{
            \makecell[{{p{\linewidth}}}]{
                \texttt{\tiny{[GM$|$GM]}}
                {'word': 'суд', 'assignment': 'team'}
            }
        }
    }
    & & \\ \\

    \theutterance \stepcounter{utterance}  
    & & & \multicolumn{2}{p{0.3\linewidth}}{
        \cellcolor[rgb]{0.9,0.9,0.9}{
            \makecell[{{p{\linewidth}}}]{
                \texttt{\tiny{[GM$|$GM]}}
                {'word': 'труба', 'assignment': 'opponent'}
            }
        }
    }
    & & \\ \\

    \theutterance \stepcounter{utterance}  
    & & & \multicolumn{2}{p{0.3\linewidth}}{
        \cellcolor[rgb]{0.9,0.9,0.9}{
            \makecell[{{p{\linewidth}}}]{
                \texttt{\tiny{[GM$|$GM]}}
                {'hidden': {'team': ['брюки', 'смотреть', 'солома', 'шип', 'бумага', 'флейта', 'замок'], 'innocent': ['лук', 'крест', 'стадион', 'крот', 'экран', 'бассейн', 'треугольник', 'пластик'], 'opponent': ['акции', 'ртуть', 'пират', 'звук', 'ветер', 'мяч', 'переключатель'], 'assassin': []}, 'revealed': {'team': {'team': ['адвокат', 'суд'], 'innocent': [], 'opponent': [], 'assassin': []}, 'opponent': {'team': [], 'innocent': [], 'opponent': ['труба'], 'assassin': []}}}
            }
        }
    }
    & & \\ \\

    \theutterance \stepcounter{utterance}  
    & & \multicolumn{4}{p{0.6\linewidth}}{
        \cellcolor[rgb]{0.9,0.9,0.9}{
            \makecell[{{p{\linewidth}}}]{
                \texttt{\tiny{[P1$\langle$GM]}}
                \texttt{Слова адвокат, суд были угаданы правильно. Ход вашего товарища по команде закончился.} \\
\texttt{Теперь предложите ещё одну подсказку, связанную с некоторыми из оставшихся слов вашей команды, и список связанных целевых слов. Помните: начните вашу подсказку с "ПОДСКАЗКА: ", поставьте новую строку и начните ваш список целевых слов, разделённых запятыми, с "ЦЕЛИ: ". Примечание: по сравнению с предыдущими запросами некоторые слова были удалены из списков.} \\
\\ 
\texttt{Оставшиеся слова вашей команды: брюки, смотреть, солома, шип, бумага, флейта, замок.} \\
\texttt{Оставшиеся слова вашего соперника: акции, ртуть, пират, труба, звук, ветер, мяч, переключатель.} \\
\texttt{Оставшиеся слова{-}отвлекатели: лук, крест, стадион, крот, экран, бассейн, треугольник, пластик.} \\
\texttt{Оставшиеся слова{-}убийцы: .} \\
            }
        }
    }
    & \\ \\

    \theutterance \stepcounter{utterance}  
    & \multicolumn{4}{p{0.6\linewidth}}{
        \cellcolor[rgb]{0.9,0.9,0.9}{
            \makecell[{{p{\linewidth}}}]{
                \texttt{\tiny{[P1$\rangle$GM]}}
                \texttt{ПОДСКАЗКА: крепость} \\
\texttt{ЦЕЛИ: замок} \\
            }
        }
    }
    & & \\ \\

    \theutterance \stepcounter{utterance}  
    & & & \multicolumn{2}{p{0.3\linewidth}}{
        \cellcolor[rgb]{0.9,0.9,0.9}{
            \makecell[{{p{\linewidth}}}]{
                \texttt{\tiny{[GM$|$GM]}}
                \texttt{крепость} \\
            }
        }
    }
    & & \\ \\

    \theutterance \stepcounter{utterance}  
    & & & \multicolumn{2}{p{0.3\linewidth}}{
        \cellcolor[rgb]{0.9,0.9,0.9}{
            \makecell[{{p{\linewidth}}}]{
                \texttt{\tiny{[GM$|$GM]}}
                ['замок']
            }
        }
    }
    & & \\ \\

    \theutterance \stepcounter{utterance}  
    & & & \multicolumn{2}{p{0.3\linewidth}}{
        \cellcolor[rgb]{0.9,0.9,0.9}{
            \makecell[{{p{\linewidth}}}]{
                \texttt{\tiny{[GM$|$GM]}}
                {'word': 'замок', 'assignment': 'team'}
            }
        }
    }
    & & \\ \\

    \theutterance \stepcounter{utterance}  
    & & \multicolumn{4}{p{0.6\linewidth}}{
        \cellcolor[rgb]{0.9,0.9,0.9}{
            \makecell[{{p{\linewidth}}}]{
                \texttt{\tiny{[GM$\rangle$P2]}}
                \texttt{Слова адвокат, суд были угаданы правильно. Ваш ход закончился.} \\
\texttt{Теперь предоставьте ещё один список, разделённый запятыми, содержащий как минимум 1 и не более чем 1 слов из следующего списка слов, которые лучше всего относятся или наиболее тесно связаны со словом 'крепость'. Помните, что нужно начать ответ с ОТВЕТ: . Обратите внимание: некоторые слова были удалены из списка по сравнению с предыдущими запросами.} \\
\\ 
\texttt{замок, ветер, шип, пират, крот, переключатель, бассейн, звук, акции, экран, брюки, пластик, солома, ртуть, мяч, смотреть, крест, треугольник, лук, бумага, стадион, флейта} \\
            }
        }
    }
    & \\ \\

    \theutterance \stepcounter{utterance}  
    & & & \multicolumn{4}{p{0.6\linewidth}}{
        \cellcolor[rgb]{0.9,0.9,0.9}{
            \makecell[{{p{\linewidth}}}]{
                \texttt{\tiny{[GM$\langle$P2]}}
                \texttt{ОТВЕТ: замок} \\
            }
        }
    }
     \\ \\

    \theutterance \stepcounter{utterance}  
    & & & \multicolumn{2}{p{0.3\linewidth}}{
        \cellcolor[rgb]{0.9,0.9,0.9}{
            \makecell[{{p{\linewidth}}}]{
                \texttt{\tiny{[GM$|$GM]}}
                ['замок']
            }
        }
    }
    & & \\ \\

    \theutterance \stepcounter{utterance}  
    & & & \multicolumn{2}{p{0.3\linewidth}}{
        \cellcolor[rgb]{0.9,0.9,0.9}{
            \makecell[{{p{\linewidth}}}]{
                \texttt{\tiny{[GM$|$GM]}}
                {'word': 'замок', 'assignment': 'team'}
            }
        }
    }
    & & \\ \\

    \theutterance \stepcounter{utterance}  
    & & & \multicolumn{2}{p{0.3\linewidth}}{
        \cellcolor[rgb]{0.9,0.9,0.9}{
            \makecell[{{p{\linewidth}}}]{
                \texttt{\tiny{[GM$|$GM]}}
                {'word': 'замок', 'assignment': 'team'}
            }
        }
    }
    & & \\ \\

    \theutterance \stepcounter{utterance}  
    & & & \multicolumn{2}{p{0.3\linewidth}}{
        \cellcolor[rgb]{0.9,0.9,0.9}{
            \makecell[{{p{\linewidth}}}]{
                \texttt{\tiny{[GM$|$GM]}}
                {'word': 'ртуть', 'assignment': 'opponent'}
            }
        }
    }
    & & \\ \\

    \theutterance \stepcounter{utterance}  
    & & & \multicolumn{2}{p{0.3\linewidth}}{
        \cellcolor[rgb]{0.9,0.9,0.9}{
            \makecell[{{p{\linewidth}}}]{
                \texttt{\tiny{[GM$|$GM]}}
                {'hidden': {'team': ['брюки', 'смотреть', 'солома', 'шип', 'бумага', 'флейта'], 'innocent': ['лук', 'крест', 'стадион', 'крот', 'экран', 'бассейн', 'треугольник', 'пластик'], 'opponent': ['акции', 'пират', 'звук', 'ветер', 'мяч', 'переключатель'], 'assassin': []}, 'revealed': {'team': {'team': ['адвокат', 'суд', 'замок'], 'innocent': [], 'opponent': [], 'assassin': []}, 'opponent': {'team': [], 'innocent': [], 'opponent': ['труба', 'ртуть'], 'assassin': []}}}
            }
        }
    }
    & & \\ \\

    \theutterance \stepcounter{utterance}  
    & & \multicolumn{4}{p{0.6\linewidth}}{
        \cellcolor[rgb]{0.9,0.9,0.9}{
            \makecell[{{p{\linewidth}}}]{
                \texttt{\tiny{[P1$\langle$GM]}}
                \texttt{Слово замок было угадано правильно. Ход вашего товарища по команде закончился.} \\
\texttt{Теперь предложите ещё одну подсказку, связанную с некоторыми из оставшихся слов вашей команды, и список связанных целевых слов. Помните: начните вашу подсказку с "ПОДСКАЗКА: ", поставьте новую строку и начните ваш список целевых слов, разделённых запятыми, с "ЦЕЛИ: ". Примечание: по сравнению с предыдущими запросами некоторые слова были удалены из списков.} \\
\\ 
\texttt{Оставшиеся слова вашей команды: брюки, смотреть, солома, шип, бумага, флейта.} \\
\texttt{Оставшиеся слова вашего соперника: акции, ртуть, пират, звук, ветер, мяч, переключатель.} \\
\texttt{Оставшиеся слова{-}отвлекатели: лук, крест, стадион, крот, экран, бассейн, треугольник, пластик.} \\
\texttt{Оставшиеся слова{-}убийцы: .} \\
            }
        }
    }
    & \\ \\

    \theutterance \stepcounter{utterance}  
    & \multicolumn{4}{p{0.6\linewidth}}{
        \cellcolor[rgb]{0.9,0.9,0.9}{
            \makecell[{{p{\linewidth}}}]{
                \texttt{\tiny{[P1$\rangle$GM]}}
                \texttt{ПОДСКАЗКА: складка} \\
\texttt{ЦЕЛИ: брюки, бумага} \\
            }
        }
    }
    & & \\ \\

    \theutterance \stepcounter{utterance}  
    & & & \multicolumn{2}{p{0.3\linewidth}}{
        \cellcolor[rgb]{0.9,0.9,0.9}{
            \makecell[{{p{\linewidth}}}]{
                \texttt{\tiny{[GM$|$GM]}}
                \texttt{складка} \\
            }
        }
    }
    & & \\ \\

    \theutterance \stepcounter{utterance}  
    & & & \multicolumn{2}{p{0.3\linewidth}}{
        \cellcolor[rgb]{0.9,0.9,0.9}{
            \makecell[{{p{\linewidth}}}]{
                \texttt{\tiny{[GM$|$GM]}}
                ['брюки', 'бумага']
            }
        }
    }
    & & \\ \\

    \theutterance \stepcounter{utterance}  
    & & & \multicolumn{2}{p{0.3\linewidth}}{
        \cellcolor[rgb]{0.9,0.9,0.9}{
            \makecell[{{p{\linewidth}}}]{
                \texttt{\tiny{[GM$|$GM]}}
                {'word': 'брюки', 'assignment': 'team'}
            }
        }
    }
    & & \\ \\

    \theutterance \stepcounter{utterance}  
    & & & \multicolumn{2}{p{0.3\linewidth}}{
        \cellcolor[rgb]{0.9,0.9,0.9}{
            \makecell[{{p{\linewidth}}}]{
                \texttt{\tiny{[GM$|$GM]}}
                {'word': 'бумага', 'assignment': 'team'}
            }
        }
    }
    & & \\ \\

    \theutterance \stepcounter{utterance}  
    & & \multicolumn{4}{p{0.6\linewidth}}{
        \cellcolor[rgb]{0.9,0.9,0.9}{
            \makecell[{{p{\linewidth}}}]{
                \texttt{\tiny{[GM$\rangle$P2]}}
                \texttt{Слово замок было угадано правильно. Ваш ход закончился.} \\
\texttt{Теперь предоставьте ещё один список, разделённый запятыми, содержащий как минимум 1 и не более чем 2 слов из следующего списка слов, которые лучше всего относятся или наиболее тесно связаны со словом 'складка'. Помните, что нужно начать ответ с ОТВЕТ: . Обратите внимание: некоторые слова были удалены из списка по сравнению с предыдущими запросами.} \\
\\ 
\texttt{ветер, шип, пират, крот, переключатель, бассейн, звук, акции, экран, брюки, пластик, солома, мяч, смотреть, крест, треугольник, лук, бумага, стадион, флейта} \\
            }
        }
    }
    & \\ \\

    \theutterance \stepcounter{utterance}  
    & & & \multicolumn{4}{p{0.6\linewidth}}{
        \cellcolor[rgb]{0.9,0.9,0.9}{
            \makecell[{{p{\linewidth}}}]{
                \texttt{\tiny{[GM$\langle$P2]}}
                \texttt{ОТВЕТ: бумага, брюки} \\
            }
        }
    }
     \\ \\

    \theutterance \stepcounter{utterance}  
    & & & \multicolumn{2}{p{0.3\linewidth}}{
        \cellcolor[rgb]{0.9,0.9,0.9}{
            \makecell[{{p{\linewidth}}}]{
                \texttt{\tiny{[GM$|$GM]}}
                ['бумага', 'брюки']
            }
        }
    }
    & & \\ \\

    \theutterance \stepcounter{utterance}  
    & & & \multicolumn{2}{p{0.3\linewidth}}{
        \cellcolor[rgb]{0.9,0.9,0.9}{
            \makecell[{{p{\linewidth}}}]{
                \texttt{\tiny{[GM$|$GM]}}
                {'word': 'бумага', 'assignment': 'team'}
            }
        }
    }
    & & \\ \\

    \theutterance \stepcounter{utterance}  
    & & & \multicolumn{2}{p{0.3\linewidth}}{
        \cellcolor[rgb]{0.9,0.9,0.9}{
            \makecell[{{p{\linewidth}}}]{
                \texttt{\tiny{[GM$|$GM]}}
                {'word': 'бумага', 'assignment': 'team'}
            }
        }
    }
    & & \\ \\

    \theutterance \stepcounter{utterance}  
    & & & \multicolumn{2}{p{0.3\linewidth}}{
        \cellcolor[rgb]{0.9,0.9,0.9}{
            \makecell[{{p{\linewidth}}}]{
                \texttt{\tiny{[GM$|$GM]}}
                {'word': 'брюки', 'assignment': 'team'}
            }
        }
    }
    & & \\ \\

    \theutterance \stepcounter{utterance}  
    & & & \multicolumn{2}{p{0.3\linewidth}}{
        \cellcolor[rgb]{0.9,0.9,0.9}{
            \makecell[{{p{\linewidth}}}]{
                \texttt{\tiny{[GM$|$GM]}}
                {'word': 'брюки', 'assignment': 'team'}
            }
        }
    }
    & & \\ \\

    \theutterance \stepcounter{utterance}  
    & & & \multicolumn{2}{p{0.3\linewidth}}{
        \cellcolor[rgb]{0.9,0.9,0.9}{
            \makecell[{{p{\linewidth}}}]{
                \texttt{\tiny{[GM$|$GM]}}
                {'word': 'пират', 'assignment': 'opponent'}
            }
        }
    }
    & & \\ \\

    \theutterance \stepcounter{utterance}  
    & & & \multicolumn{2}{p{0.3\linewidth}}{
        \cellcolor[rgb]{0.9,0.9,0.9}{
            \makecell[{{p{\linewidth}}}]{
                \texttt{\tiny{[GM$|$GM]}}
                {'hidden': {'team': ['смотреть', 'солома', 'шип', 'флейта'], 'innocent': ['лук', 'крест', 'стадион', 'крот', 'экран', 'бассейн', 'треугольник', 'пластик'], 'opponent': ['акции', 'звук', 'ветер', 'мяч', 'переключатель'], 'assassin': []}, 'revealed': {'team': {'team': ['адвокат', 'суд', 'замок', 'бумага', 'брюки'], 'innocent': [], 'opponent': [], 'assassin': []}, 'opponent': {'team': [], 'innocent': [], 'opponent': ['труба', 'ртуть', 'пират'], 'assassin': []}}}
            }
        }
    }
    & & \\ \\

    \theutterance \stepcounter{utterance}  
    & & \multicolumn{4}{p{0.6\linewidth}}{
        \cellcolor[rgb]{0.9,0.9,0.9}{
            \makecell[{{p{\linewidth}}}]{
                \texttt{\tiny{[P1$\langle$GM]}}
                \texttt{Слова бумага, брюки были угаданы правильно. Ход вашего товарища по команде закончился.} \\
\texttt{Теперь предложите ещё одну подсказку, связанную с некоторыми из оставшихся слов вашей команды, и список связанных целевых слов. Помните: начните вашу подсказку с "ПОДСКАЗКА: ", поставьте новую строку и начните ваш список целевых слов, разделённых запятыми, с "ЦЕЛИ: ". Примечание: по сравнению с предыдущими запросами некоторые слова были удалены из списков.} \\
\\ 
\texttt{Оставшиеся слова вашей команды: смотреть, солома, шип, флейта.} \\
\texttt{Оставшиеся слова вашего соперника: акции, пират, звук, ветер, мяч, переключатель.} \\
\texttt{Оставшиеся слова{-}отвлекатели: лук, крест, стадион, крот, экран, бассейн, треугольник, пластик.} \\
\texttt{Оставшиеся слова{-}убийцы: .} \\
            }
        }
    }
    & \\ \\

    \theutterance \stepcounter{utterance}  
    & \multicolumn{4}{p{0.6\linewidth}}{
        \cellcolor[rgb]{0.9,0.9,0.9}{
            \makecell[{{p{\linewidth}}}]{
                \texttt{\tiny{[P1$\rangle$GM]}}
                \texttt{ПОДСКАЗКА: взгляд} \\
\texttt{ЦЕЛИ: смотреть} \\
            }
        }
    }
    & & \\ \\

    \theutterance \stepcounter{utterance}  
    & & & \multicolumn{2}{p{0.3\linewidth}}{
        \cellcolor[rgb]{0.9,0.9,0.9}{
            \makecell[{{p{\linewidth}}}]{
                \texttt{\tiny{[GM$|$GM]}}
                \texttt{взгляд} \\
            }
        }
    }
    & & \\ \\

    \theutterance \stepcounter{utterance}  
    & & & \multicolumn{2}{p{0.3\linewidth}}{
        \cellcolor[rgb]{0.9,0.9,0.9}{
            \makecell[{{p{\linewidth}}}]{
                \texttt{\tiny{[GM$|$GM]}}
                ['смотреть']
            }
        }
    }
    & & \\ \\

    \theutterance \stepcounter{utterance}  
    & & & \multicolumn{2}{p{0.3\linewidth}}{
        \cellcolor[rgb]{0.9,0.9,0.9}{
            \makecell[{{p{\linewidth}}}]{
                \texttt{\tiny{[GM$|$GM]}}
                {'word': 'смотреть', 'assignment': 'team'}
            }
        }
    }
    & & \\ \\

    \theutterance \stepcounter{utterance}  
    & & \multicolumn{4}{p{0.6\linewidth}}{
        \cellcolor[rgb]{0.9,0.9,0.9}{
            \makecell[{{p{\linewidth}}}]{
                \texttt{\tiny{[GM$\rangle$P2]}}
                \texttt{Слова бумага, брюки были угаданы правильно. Ваш ход закончился.} \\
\texttt{Теперь предоставьте ещё один список, разделённый запятыми, содержащий как минимум 1 и не более чем 1 слов из следующего списка слов, которые лучше всего относятся или наиболее тесно связаны со словом 'взгляд'. Помните, что нужно начать ответ с ОТВЕТ: . Обратите внимание: некоторые слова были удалены из списка по сравнению с предыдущими запросами.} \\
\\ 
\texttt{ветер, шип, крот, переключатель, бассейн, звук, акции, экран, пластик, солома, мяч, смотреть, крест, треугольник, лук, стадион, флейта} \\
            }
        }
    }
    & \\ \\

    \theutterance \stepcounter{utterance}  
    & & & \multicolumn{4}{p{0.6\linewidth}}{
        \cellcolor[rgb]{0.9,0.9,0.9}{
            \makecell[{{p{\linewidth}}}]{
                \texttt{\tiny{[GM$\langle$P2]}}
                \texttt{ОТВЕТ: смотреть} \\
            }
        }
    }
     \\ \\

    \theutterance \stepcounter{utterance}  
    & & & \multicolumn{2}{p{0.3\linewidth}}{
        \cellcolor[rgb]{0.9,0.9,0.9}{
            \makecell[{{p{\linewidth}}}]{
                \texttt{\tiny{[GM$|$GM]}}
                ['смотреть']
            }
        }
    }
    & & \\ \\

    \theutterance \stepcounter{utterance}  
    & & & \multicolumn{2}{p{0.3\linewidth}}{
        \cellcolor[rgb]{0.9,0.9,0.9}{
            \makecell[{{p{\linewidth}}}]{
                \texttt{\tiny{[GM$|$GM]}}
                {'word': 'смотреть', 'assignment': 'team'}
            }
        }
    }
    & & \\ \\

    \theutterance \stepcounter{utterance}  
    & & & \multicolumn{2}{p{0.3\linewidth}}{
        \cellcolor[rgb]{0.9,0.9,0.9}{
            \makecell[{{p{\linewidth}}}]{
                \texttt{\tiny{[GM$|$GM]}}
                {'word': 'смотреть', 'assignment': 'team'}
            }
        }
    }
    & & \\ \\

    \theutterance \stepcounter{utterance}  
    & & & \multicolumn{2}{p{0.3\linewidth}}{
        \cellcolor[rgb]{0.9,0.9,0.9}{
            \makecell[{{p{\linewidth}}}]{
                \texttt{\tiny{[GM$|$GM]}}
                {'word': 'акции', 'assignment': 'opponent'}
            }
        }
    }
    & & \\ \\

    \theutterance \stepcounter{utterance}  
    & & & \multicolumn{2}{p{0.3\linewidth}}{
        \cellcolor[rgb]{0.9,0.9,0.9}{
            \makecell[{{p{\linewidth}}}]{
                \texttt{\tiny{[GM$|$GM]}}
                {'hidden': {'team': ['солома', 'шип', 'флейта'], 'innocent': ['лук', 'крест', 'стадион', 'крот', 'экран', 'бассейн', 'треугольник', 'пластик'], 'opponent': ['звук', 'ветер', 'мяч', 'переключатель'], 'assassin': []}, 'revealed': {'team': {'team': ['адвокат', 'суд', 'замок', 'бумага', 'брюки', 'смотреть'], 'innocent': [], 'opponent': [], 'assassin': []}, 'opponent': {'team': [], 'innocent': [], 'opponent': ['труба', 'ртуть', 'пират', 'акции'], 'assassin': []}}}
            }
        }
    }
    & & \\ \\

    \theutterance \stepcounter{utterance}  
    & & \multicolumn{4}{p{0.6\linewidth}}{
        \cellcolor[rgb]{0.9,0.9,0.9}{
            \makecell[{{p{\linewidth}}}]{
                \texttt{\tiny{[P1$\langle$GM]}}
                \texttt{Слово смотреть было угадано правильно. Ход вашего товарища по команде закончился.} \\
\texttt{Теперь предложите ещё одну подсказку, связанную с некоторыми из оставшихся слов вашей команды, и список связанных целевых слов. Помните: начните вашу подсказку с "ПОДСКАЗКА: ", поставьте новую строку и начните ваш список целевых слов, разделённых запятыми, с "ЦЕЛИ: ". Примечание: по сравнению с предыдущими запросами некоторые слова были удалены из списков.} \\
\\ 
\texttt{Оставшиеся слова вашей команды: солома, шип, флейта.} \\
\texttt{Оставшиеся слова вашего соперника: акции, звук, ветер, мяч, переключатель.} \\
\texttt{Оставшиеся слова{-}отвлекатели: лук, крест, стадион, крот, экран, бассейн, треугольник, пластик.} \\
\texttt{Оставшиеся слова{-}убийцы: .} \\
            }
        }
    }
    & \\ \\

    \theutterance \stepcounter{utterance}  
    & \multicolumn{4}{p{0.6\linewidth}}{
        \cellcolor[rgb]{0.9,0.9,0.9}{
            \makecell[{{p{\linewidth}}}]{
                \texttt{\tiny{[P1$\rangle$GM]}}
                \texttt{ПОДСКАЗКА: колючка} \\
\texttt{ЦЕЛИ: шип} \\
            }
        }
    }
    & & \\ \\

    \theutterance \stepcounter{utterance}  
    & & & \multicolumn{2}{p{0.3\linewidth}}{
        \cellcolor[rgb]{0.9,0.9,0.9}{
            \makecell[{{p{\linewidth}}}]{
                \texttt{\tiny{[GM$|$GM]}}
                \texttt{колючка} \\
            }
        }
    }
    & & \\ \\

    \theutterance \stepcounter{utterance}  
    & & & \multicolumn{2}{p{0.3\linewidth}}{
        \cellcolor[rgb]{0.9,0.9,0.9}{
            \makecell[{{p{\linewidth}}}]{
                \texttt{\tiny{[GM$|$GM]}}
                ['шип']
            }
        }
    }
    & & \\ \\

    \theutterance \stepcounter{utterance}  
    & & & \multicolumn{2}{p{0.3\linewidth}}{
        \cellcolor[rgb]{0.9,0.9,0.9}{
            \makecell[{{p{\linewidth}}}]{
                \texttt{\tiny{[GM$|$GM]}}
                {'word': 'шип', 'assignment': 'team'}
            }
        }
    }
    & & \\ \\

    \theutterance \stepcounter{utterance}  
    & & \multicolumn{4}{p{0.6\linewidth}}{
        \cellcolor[rgb]{0.9,0.9,0.9}{
            \makecell[{{p{\linewidth}}}]{
                \texttt{\tiny{[GM$\rangle$P2]}}
                \texttt{Слово смотреть было угадано правильно. Ваш ход закончился.} \\
\texttt{Теперь предоставьте ещё один список, разделённый запятыми, содержащий как минимум 1 и не более чем 1 слов из следующего списка слов, которые лучше всего относятся или наиболее тесно связаны со словом 'колючка'. Помните, что нужно начать ответ с ОТВЕТ: . Обратите внимание: некоторые слова были удалены из списка по сравнению с предыдущими запросами.} \\
\\ 
\texttt{ветер, шип, крот, переключатель, бассейн, звук, экран, пластик, солома, мяч, крест, треугольник, лук, стадион, флейта} \\
            }
        }
    }
    & \\ \\

    \theutterance \stepcounter{utterance}  
    & & & \multicolumn{4}{p{0.6\linewidth}}{
        \cellcolor[rgb]{0.9,0.9,0.9}{
            \makecell[{{p{\linewidth}}}]{
                \texttt{\tiny{[GM$\langle$P2]}}
                \texttt{ОТВЕТ: шип} \\
            }
        }
    }
     \\ \\

    \theutterance \stepcounter{utterance}  
    & & & \multicolumn{2}{p{0.3\linewidth}}{
        \cellcolor[rgb]{0.9,0.9,0.9}{
            \makecell[{{p{\linewidth}}}]{
                \texttt{\tiny{[GM$|$GM]}}
                ['шип']
            }
        }
    }
    & & \\ \\

    \theutterance \stepcounter{utterance}  
    & & & \multicolumn{2}{p{0.3\linewidth}}{
        \cellcolor[rgb]{0.9,0.9,0.9}{
            \makecell[{{p{\linewidth}}}]{
                \texttt{\tiny{[GM$|$GM]}}
                {'word': 'шип', 'assignment': 'team'}
            }
        }
    }
    & & \\ \\

    \theutterance \stepcounter{utterance}  
    & & & \multicolumn{2}{p{0.3\linewidth}}{
        \cellcolor[rgb]{0.9,0.9,0.9}{
            \makecell[{{p{\linewidth}}}]{
                \texttt{\tiny{[GM$|$GM]}}
                {'word': 'шип', 'assignment': 'team'}
            }
        }
    }
    & & \\ \\

    \theutterance \stepcounter{utterance}  
    & & & \multicolumn{2}{p{0.3\linewidth}}{
        \cellcolor[rgb]{0.9,0.9,0.9}{
            \makecell[{{p{\linewidth}}}]{
                \texttt{\tiny{[GM$|$GM]}}
                {'word': 'звук', 'assignment': 'opponent'}
            }
        }
    }
    & & \\ \\

    \theutterance \stepcounter{utterance}  
    & & & \multicolumn{2}{p{0.3\linewidth}}{
        \cellcolor[rgb]{0.9,0.9,0.9}{
            \makecell[{{p{\linewidth}}}]{
                \texttt{\tiny{[GM$|$GM]}}
                {'hidden': {'team': ['солома', 'флейта'], 'innocent': ['лук', 'крест', 'стадион', 'крот', 'экран', 'бассейн', 'треугольник', 'пластик'], 'opponent': ['ветер', 'мяч', 'переключатель'], 'assassin': []}, 'revealed': {'team': {'team': ['адвокат', 'суд', 'замок', 'бумага', 'брюки', 'смотреть', 'шип'], 'innocent': [], 'opponent': [], 'assassin': []}, 'opponent': {'team': [], 'innocent': [], 'opponent': ['труба', 'ртуть', 'пират', 'акции', 'звук'], 'assassin': []}}}
            }
        }
    }
    & & \\ \\

    \theutterance \stepcounter{utterance}  
    & & \multicolumn{4}{p{0.6\linewidth}}{
        \cellcolor[rgb]{0.9,0.9,0.9}{
            \makecell[{{p{\linewidth}}}]{
                \texttt{\tiny{[P1$\langle$GM]}}
                \texttt{Слово шип было угадано правильно. Ход вашего товарища по команде закончился.} \\
\texttt{Теперь предложите ещё одну подсказку, связанную с некоторыми из оставшихся слов вашей команды, и список связанных целевых слов. Помните: начните вашу подсказку с "ПОДСКАЗКА: ", поставьте новую строку и начните ваш список целевых слов, разделённых запятыми, с "ЦЕЛИ: ". Примечание: по сравнению с предыдущими запросами некоторые слова были удалены из списков.} \\
\\ 
\texttt{Оставшиеся слова вашей команды: солома, флейта.} \\
\texttt{Оставшиеся слова вашего соперника: звук, ветер, мяч, переключатель.} \\
\texttt{Оставшиеся слова{-}отвлекатели: лук, крест, стадион, крот, экран, бассейн, треугольник, пластик.} \\
\texttt{Оставшиеся слова{-}убийцы: .} \\
            }
        }
    }
    & \\ \\

    \theutterance \stepcounter{utterance}  
    & \multicolumn{4}{p{0.6\linewidth}}{
        \cellcolor[rgb]{0.9,0.9,0.9}{
            \makecell[{{p{\linewidth}}}]{
                \texttt{\tiny{[P1$\rangle$GM]}}
                \texttt{ПОДСКАЗКА: сено} \\
\texttt{ЦЕЛИ: солома} \\
            }
        }
    }
    & & \\ \\

    \theutterance \stepcounter{utterance}  
    & & & \multicolumn{2}{p{0.3\linewidth}}{
        \cellcolor[rgb]{0.9,0.9,0.9}{
            \makecell[{{p{\linewidth}}}]{
                \texttt{\tiny{[GM$|$GM]}}
                \texttt{сено} \\
            }
        }
    }
    & & \\ \\

    \theutterance \stepcounter{utterance}  
    & & & \multicolumn{2}{p{0.3\linewidth}}{
        \cellcolor[rgb]{0.9,0.9,0.9}{
            \makecell[{{p{\linewidth}}}]{
                \texttt{\tiny{[GM$|$GM]}}
                ['солома']
            }
        }
    }
    & & \\ \\

    \theutterance \stepcounter{utterance}  
    & & & \multicolumn{2}{p{0.3\linewidth}}{
        \cellcolor[rgb]{0.9,0.9,0.9}{
            \makecell[{{p{\linewidth}}}]{
                \texttt{\tiny{[GM$|$GM]}}
                {'word': 'солома', 'assignment': 'team'}
            }
        }
    }
    & & \\ \\

    \theutterance \stepcounter{utterance}  
    & & \multicolumn{4}{p{0.6\linewidth}}{
        \cellcolor[rgb]{0.9,0.9,0.9}{
            \makecell[{{p{\linewidth}}}]{
                \texttt{\tiny{[GM$\rangle$P2]}}
                \texttt{Слово шип было угадано правильно. Ваш ход закончился.} \\
\texttt{Теперь предоставьте ещё один список, разделённый запятыми, содержащий как минимум 1 и не более чем 1 слов из следующего списка слов, которые лучше всего относятся или наиболее тесно связаны со словом 'сено'. Помните, что нужно начать ответ с ОТВЕТ: . Обратите внимание: некоторые слова были удалены из списка по сравнению с предыдущими запросами.} \\
\\ 
\texttt{ветер, крот, переключатель, бассейн, экран, пластик, солома, мяч, крест, треугольник, лук, стадион, флейта} \\
            }
        }
    }
    & \\ \\

    \theutterance \stepcounter{utterance}  
    & & & \multicolumn{4}{p{0.6\linewidth}}{
        \cellcolor[rgb]{0.9,0.9,0.9}{
            \makecell[{{p{\linewidth}}}]{
                \texttt{\tiny{[GM$\langle$P2]}}
                \texttt{ОТВЕТ: солома} \\
            }
        }
    }
     \\ \\

    \theutterance \stepcounter{utterance}  
    & & & \multicolumn{2}{p{0.3\linewidth}}{
        \cellcolor[rgb]{0.9,0.9,0.9}{
            \makecell[{{p{\linewidth}}}]{
                \texttt{\tiny{[GM$|$GM]}}
                ['солома']
            }
        }
    }
    & & \\ \\

    \theutterance \stepcounter{utterance}  
    & & & \multicolumn{2}{p{0.3\linewidth}}{
        \cellcolor[rgb]{0.9,0.9,0.9}{
            \makecell[{{p{\linewidth}}}]{
                \texttt{\tiny{[GM$|$GM]}}
                {'word': 'солома', 'assignment': 'team'}
            }
        }
    }
    & & \\ \\

    \theutterance \stepcounter{utterance}  
    & & & \multicolumn{2}{p{0.3\linewidth}}{
        \cellcolor[rgb]{0.9,0.9,0.9}{
            \makecell[{{p{\linewidth}}}]{
                \texttt{\tiny{[GM$|$GM]}}
                {'word': 'солома', 'assignment': 'team'}
            }
        }
    }
    & & \\ \\

    \theutterance \stepcounter{utterance}  
    & & & \multicolumn{2}{p{0.3\linewidth}}{
        \cellcolor[rgb]{0.9,0.9,0.9}{
            \makecell[{{p{\linewidth}}}]{
                \texttt{\tiny{[GM$|$GM]}}
                {'word': 'переключатель', 'assignment': 'opponent'}
            }
        }
    }
    & & \\ \\

    \theutterance \stepcounter{utterance}  
    & & & \multicolumn{2}{p{0.3\linewidth}}{
        \cellcolor[rgb]{0.9,0.9,0.9}{
            \makecell[{{p{\linewidth}}}]{
                \texttt{\tiny{[GM$|$GM]}}
                {'hidden': {'team': ['флейта'], 'innocent': ['лук', 'крест', 'стадион', 'крот', 'экран', 'бассейн', 'треугольник', 'пластик'], 'opponent': ['ветер', 'мяч'], 'assassin': []}, 'revealed': {'team': {'team': ['адвокат', 'суд', 'замок', 'бумага', 'брюки', 'смотреть', 'шип', 'солома'], 'innocent': [], 'opponent': [], 'assassin': []}, 'opponent': {'team': [], 'innocent': [], 'opponent': ['труба', 'ртуть', 'пират', 'акции', 'звук', 'переключатель'], 'assassin': []}}}
            }
        }
    }
    & & \\ \\

    \theutterance \stepcounter{utterance}  
    & & \multicolumn{4}{p{0.6\linewidth}}{
        \cellcolor[rgb]{0.9,0.9,0.9}{
            \makecell[{{p{\linewidth}}}]{
                \texttt{\tiny{[P1$\langle$GM]}}
                \texttt{Слово солома было угадано правильно. Ход вашего товарища по команде закончился.} \\
\texttt{Теперь предложите ещё одну подсказку, связанную с некоторыми из оставшихся слов вашей команды, и список связанных целевых слов. Помните: начните вашу подсказку с "ПОДСКАЗКА: ", поставьте новую строку и начните ваш список целевых слов, разделённых запятыми, с "ЦЕЛИ: ". Примечание: по сравнению с предыдущими запросами некоторые слова были удалены из списков.} \\
\\ 
\texttt{Оставшиеся слова вашей команды: флейта.} \\
\texttt{Оставшиеся слова вашего соперника: ветер, мяч, переключатель.} \\
\texttt{Оставшиеся слова{-}отвлекатели: лук, крест, стадион, крот, экран, бассейн, треугольник, пластик.} \\
\texttt{Оставшиеся слова{-}убийцы: .} \\
            }
        }
    }
    & \\ \\

    \theutterance \stepcounter{utterance}  
    & \multicolumn{4}{p{0.6\linewidth}}{
        \cellcolor[rgb]{0.9,0.9,0.9}{
            \makecell[{{p{\linewidth}}}]{
                \texttt{\tiny{[P1$\rangle$GM]}}
                \texttt{ПОДСКАЗКА: инструмент} \\
\texttt{ЦЕЛИ: флейта} \\
            }
        }
    }
    & & \\ \\

    \theutterance \stepcounter{utterance}  
    & & & \multicolumn{2}{p{0.3\linewidth}}{
        \cellcolor[rgb]{0.9,0.9,0.9}{
            \makecell[{{p{\linewidth}}}]{
                \texttt{\tiny{[GM$|$GM]}}
                \texttt{инструмент} \\
            }
        }
    }
    & & \\ \\

    \theutterance \stepcounter{utterance}  
    & & & \multicolumn{2}{p{0.3\linewidth}}{
        \cellcolor[rgb]{0.9,0.9,0.9}{
            \makecell[{{p{\linewidth}}}]{
                \texttt{\tiny{[GM$|$GM]}}
                ['флейта']
            }
        }
    }
    & & \\ \\

    \theutterance \stepcounter{utterance}  
    & & & \multicolumn{2}{p{0.3\linewidth}}{
        \cellcolor[rgb]{0.9,0.9,0.9}{
            \makecell[{{p{\linewidth}}}]{
                \texttt{\tiny{[GM$|$GM]}}
                {'word': 'флейта', 'assignment': 'team'}
            }
        }
    }
    & & \\ \\

    \theutterance \stepcounter{utterance}  
    & & \multicolumn{4}{p{0.6\linewidth}}{
        \cellcolor[rgb]{0.9,0.9,0.9}{
            \makecell[{{p{\linewidth}}}]{
                \texttt{\tiny{[GM$\rangle$P2]}}
                \texttt{Слово солома было угадано правильно. Ваш ход закончился.} \\
\texttt{Теперь предоставьте ещё один список, разделённый запятыми, содержащий как минимум 1 и не более чем 1 слов из следующего списка слов, которые лучше всего относятся или наиболее тесно связаны со словом 'инструмент'. Помните, что нужно начать ответ с ОТВЕТ: . Обратите внимание: некоторые слова были удалены из списка по сравнению с предыдущими запросами.} \\
\\ 
\texttt{ветер, крот, бассейн, экран, пластик, мяч, крест, треугольник, лук, стадион, флейта} \\
            }
        }
    }
    & \\ \\

    \theutterance \stepcounter{utterance}  
    & & & \multicolumn{4}{p{0.6\linewidth}}{
        \cellcolor[rgb]{0.9,0.9,0.9}{
            \makecell[{{p{\linewidth}}}]{
                \texttt{\tiny{[GM$\langle$P2]}}
                \texttt{ОТВЕТ: флейта} \\
            }
        }
    }
     \\ \\

    \theutterance \stepcounter{utterance}  
    & & & \multicolumn{2}{p{0.3\linewidth}}{
        \cellcolor[rgb]{0.9,0.9,0.9}{
            \makecell[{{p{\linewidth}}}]{
                \texttt{\tiny{[GM$|$GM]}}
                ['флейта']
            }
        }
    }
    & & \\ \\

    \theutterance \stepcounter{utterance}  
    & & & \multicolumn{2}{p{0.3\linewidth}}{
        \cellcolor[rgb]{0.9,0.9,0.9}{
            \makecell[{{p{\linewidth}}}]{
                \texttt{\tiny{[GM$|$GM]}}
                {'word': 'флейта', 'assignment': 'team'}
            }
        }
    }
    & & \\ \\

    \theutterance \stepcounter{utterance}  
    & & & \multicolumn{2}{p{0.3\linewidth}}{
        \cellcolor[rgb]{0.9,0.9,0.9}{
            \makecell[{{p{\linewidth}}}]{
                \texttt{\tiny{[GM$|$GM]}}
                {'word': 'флейта', 'assignment': 'team'}
            }
        }
    }
    & & \\ \\

    \theutterance \stepcounter{utterance}  
    & & & \multicolumn{2}{p{0.3\linewidth}}{
        \cellcolor[rgb]{0.9,0.9,0.9}{
            \makecell[{{p{\linewidth}}}]{
                \texttt{\tiny{[GM$|$GM]}}
                \texttt{team has won} \\
            }
        }
    }
    & & \\ \\

\end{supertabular}
}

\end{document}
